\section{引言}
\label{sec:introduction}

以神经特征构造数值试验空间，为偏微分方程提供了多项式基之外的选择。然而，一个流动求解器同时包含函数表示、空间投影、时间积分、线性代数、非线性变换与运行时管理。将这些部分一起替换之后获得的加速，不能自动归于随机特征。本研究针对这一归因问题，考察固定随机特征 Galerkin 方法在二维非稳态不可压流中的数值可行性，以及其数学结构如何转化为可测量的计算收益。

非稳态尾迹使比较尤其困难。速度是流函数的一阶导数，涡量涉及二阶导数；速度误差较小的近似，仍可能给出较差的涡量。交替涡脱落还会积累相位误差。如果允许每条轨迹独立平移时间，图像可能相似而原始初边值问题的解已不同。另一方面，在小型 GPU 问题上，一次算子调用只有几十微秒，设备编译、模块装载、FFT 计划及上下文初始化却可能占据可观的完整求解时间。因此，本文把共同物理精度、原始时间坐标以及端到端成本同时作为比较条件。

研究对象是有限圆环内、无反射约束的 \(Re=100\) 流动，积分至 \(T=72\)。径向分别采用固定随机三角特征、确定性 Legendre 特征与四阶空间有限差分；角向保留完整实 Fourier 表示。三个方法族使用一致的物理问题及二阶时间积分。Legendre--Galerkin 本身就是传统 CFD 对照，不应因其使用谱基而被排除。本文并不将固定特征 Galerkin 求解器等同于通过优化器训练整段轨迹的 PINN。

本文围绕以下三个问题组织方法与证据：第一，在保留速度与涡量精度时，分段时间步和较小试验空间能够减少多少实际工作？第二，共同相位坐标、FFT 融合和并行观测分别改变哪些计算环节，代数简化是否仍被实现中的其他成本抵消？第三，在相同精度、初始化政策和计时边界下，当前证据是否足以支持相对强确定性对照的整体速度优势？这三个问题分别需要离散分析、算子测试和完整隔离实验，不能由单一加速比回答。

\begin{table}[!htbp]
\centering\small
\caption{本文工作的归属与可支持的结论。数学适配不等于全新基本算法，局部实现收益不等于完整求解优势。}
\label{tab:contribution-map}
\begin{tabularx}{\linewidth}{@{}>{\raggedright\arraybackslash}p{.20\linewidth}XXX@{}}
\toprule
工作 & 类别 & 已有证据 & 结论边界 \\
\midrule
边界提升与一次变步桥接 & 数值方法适配 & 局部一致性推导及整轨迹资格 & 无完整 NS 无条件稳定性定理 \\
共同相位坐标与锚点 & 线性代数分析 & 同角模代数恒等式及上界 & 上界最优不等于实测耗时最优 \\
径向保留空间与 Legendre 对照 & 离散设计与验证 & 五个固定种子及确定性降秩对照 & 无随机特征普适优势证明 \\
融合 FFT、warp 观测、C 内核 & 底层算子工程 & 源码、局部与短程验证 & 不属于第三方 FFT 库内部补丁 \\
准备复用、缓存与启动路径 & 运行工程 & 分阶段记录与缓存等价性 & 最终隔离排名尚未建立 \\
\bottomrule
\end{tabularx}
\end{table}

本文的主要结果是可复核的工作量减少和性能归因，而非普遍速度胜出。两档时间策略将真实非线性求值减少约 \(11.11\%\) 和 \(9.36\%\)；最终 39 个配置在两种缓存政策下产生的 78 条 GPU 轨迹全部通过登记的数值门槛。专用融合 FFT 在固定状态和短程推进中有效，但首次完整求解变慢。最终完整隔离批次未满足接受条件，因此结果部分不据其残余记录生成方法排名。这一边界决定了本文应被理解为数值计算与应用高性能计算研究，现有理论和实验尚不足以构成一般收敛理论或普适优于传统 CFD 的论断。

补充的两因素消融进一步检验随机性的必要性：固定频率中点且保留解析层函数时，两档配置仍通过共同场误差门槛；删除非零模解析层函数则使本次12个配置均不合格，但同时降低了保留秩。本文因此把结论限定为当前富集空间的有效性，不将其自动归因于随机抽样。新增径向/模态误差分解及延后换步对照见第~\ref{sec:review-experiments}节，非可分几何下丢失的代数结构见第~\ref{sec:geometry-scope}节。

\section{相关工作与研究定位}
\label{sec:related-work}

\subsection{随机特征、投影求解与训练型方法}

随机特征方法通过固定非线性特征并确定组合系数求解 PDE\cite{r1}。局部极限学习机将固定隐藏层与区域分解结合\cite{r2}，随机神经网络 Petrov--Galerkin 方法已用于弹性及 Navier--Stokes 方程\cite{r3}。高精度 RFM、随机预条件和正交特征学习分别讨论误差、病态性及表示质量\cite{r4,r5,r6}。这些工作意味着，随机基与 Galerkin 的结合、预条件或正交化均不能单独作为本文的新颖性主张。

时间依赖 RFM 还可使用时空特征、分片统一和整体或分块最小二乘求解\cite{r15}。本文采用另一条实现路径：仅径向试验空间使用固定随机特征，角向使用 Fourier，时间用逐步 IMEX 积分。它没有学习时间特征，也没有沿整条轨迹反复最小化残差。因而针对时空最小二乘块传递的理论，不能未经推导套用于这里的 SBDF2 演化；其相反方向的性能结论也不能直接移植。

PINN 通过自动微分构造方程残差并训练网络参数\cite{r7}。因果训练调整演化问题的训练次序，PirateNets 用自适应残差结构改善深层 PINN 的训练\cite{r8,r16}。与这些方法比较时，应将训练、重启、采样、预测及导数重构按用途计入成本。已有论文中的误差和时间若采用不同几何、边界、终止时间或硬件，不能拼入本文的同精度表。当前结果没有这类同问题、同门槛的完整训练型对照，因此相关文献用于界定研究范围，不用于宣称本文更快。

\subsection{时间离散与线性代数结构}

IMEX 多步法及其变步推广已有系统研究\cite{r9,r12}。本文采用的二阶差分系数属于这一经典框架，新增分析针对当前离散质量、边界提升、历史通量包及一次步长切换的兼容关系。类似地，对称正定广义特征问题的白化与对角化是标准线性代数工具\cite{r10}。这里值得分析的是：同一角模多个时间相位能否共享坐标；如何界定锚点选择的作用；以及为何保持旧设备接口时，较少分解未必降低准备时间。

\subsection{GPU 融合与应用层性能}

FFT 去混叠、实场 Hermitian 结构和混合基算法都有成熟来源\cite{r11,r13}。cuFFTDx 已提供设备内 FFT、融合卷积及与 cuFFT 的比较示例\cite{r17}，warp 内交换和归约也属于 CUDA 标准编程机制\cite{r18}。本文的工程贡献是针对特定 FP64 多场通量布局实现并验证这些机制，保留强弱通道误差、启动负担和失败分支；不能把已有融合原理重新包装成新的 FFT 算法。

高层语言可以组织高性能流动求解，而计算密集部分由编译后的数值库或专用内核执行\cite{r14}。因此，本文按数据移动、同步、矩阵布局和生命周期讨论瓶颈。是否采用 C、CUDA 或其他宿主语言，应由所占成本和可减少的工作决定。相同原则也适用于开源贡献：本项目重写的是应用侧算子，当前没有修改第三方库源码。向上游贡献一个可复用样例，需要独立正确性、等价库实现对照和可复现构建，而不能仅凭本项目中的局部加速提出库性能修复。
